\scanpage{28}
par exemple si $T\quad S$ est\unsure{h} plat, ou dans le\deleted{s cas} kummérien\deleted{s} \inserted{a)} envisagé\deleted{s} plus haut (en supposant que les images inverses des $D_i$ dans $T$ existent, et prenant ces derniers pour définir une donnée de recollement sur $T$).

On recule devant la tâche d'expliciter une propriété de transitivi- pour un composé $U\quad T\quad S$, et jetant provisoirement un voile pudique sur ces abimes, nous allons remarquer plûtôt que le foncteur $\mathrm{Rev}^{\underline R(S)}(S)\qquad\mathrm{Rev}^{\underline R(T)}(T)$ qu'on vient de construire a le bon goût d'être \uline{exact}. Si alors , dans les conditions de l'exemple 3, $\quad$ est un point géométrique de $T$, \deleted{non ramifié pour les données $\underline R_{(\unsure{T})}$}\inserted{non ramifié pour $\underline R(T)$,} dont l'imgge dans $S$ est non ramifié pour $\underline R_{(S)}$, et si on suppose que \deleted{$\underline R_{(T)}$ est l'image inverse de $\underline R_{(S)}$, alors} l'homomorphisme \inserted{$u$ dans} $(\unsure{*|=})$ de foncteurs composés induit un isomorphisme sur les fibres géométriques en $\quad$, alors les sorites de Exp V impliquent qu'on a un homomorphisme canonique $ {}_{1}^{\underline R(\deleted{S}\unsure{T|S})}(\deleted{S}T,{}_{\deleted{S}T})\qquad{}_{1}^{\underline R(S)}(S,{}_{S})$. Ceci s'applique en particulier au cas de données kummériennes type a), ou type b) \inserted{$1^{\underline{\mathrm o}}$ et $2^{\underline{\mathrm o}}$;} quand on suppose que $T\quad S$ est un morphisme régulier \unsure{resp.|repp.} lisse (de sorte que l\deleted{es}'image\deleted{s} inverse\deleted{s} de\deleted{s} $D_{\deleted{i}}$ satisf\unsure{a}it encore aux mêmes conditions de ``normal crossings'' que $D$ lui-même).

Il y a a expliciter deux cas particuliers importants de fonctorialité. L'un est celui où $T\quad S$ est étale, et où $\underline R_{(T)}$ est induit par $\underline R_{(S)}$ sur $S$, auquel le \deleted{foncteur} carré nous était déjà connu par ailleurs (et l'homomorphisme de foncteurs composés y est un isomorphisme), puisque l'on dispose d'un foncteur cartésien $\underline{\mathrm{Rev}}^{R}\qquad\underline{\mathrm{Rev}}$ sur $\underline{\mathrm{Et}}_S$. Lorsqu'on suppose de plus $T$ fini sur $S$, $S$ et $T$ connexes, et qu'on est dans les conditions de l'exemple 3, alors l'homomorphisme sur les groupes \unsure{fondamentau|fondamentaux} est un monomorphisme, dont l'image est le sous-groupe de $ {}_{1}^{\underline R}(S,\quad)$ défini par le revêtement ponctué $T$ de $S$. Le deuxième cas est celui \unsure{où} $T\quad S$ est \unsure{un} \unsure{morphisme} entier, radiciel et surjectif, et où $\underline R_{(T)}$
% archive sheet number: 27
